\section{Related Work}

\subsection{LLM-Driven Multi-Robot Collaboration}

Several recent systems have demonstrated that LLMs can serve as high-level planners for multi-robot task allocation. COHERENT [cite] decomposes rescue missions into sub-tasks and allocates them to heterogeneous robots through a propose–execute–feedback loop, achieving strong performance on simulated disaster scenarios. RoCo [cite] and ReAd [cite] use multi-agent dialogue between LLM instances to reach consensus on robot assignments. MetaGPT [cite] formalises multi-agent LLM workflows using structured role assignments. None of these systems incorporate a human operator in the planning loop, and none use a standardised protocol layer such as MCP for tool interaction.

\subsection{Human-in-the-Loop Robot Systems}
Human oversight of robot teams has a long history in the shared-autonomy literature. Classical approaches use mixed-initiative control [cite] or sliding autonomy [cite], where the system's autonomy level is adjusted based on task difficulty or robot confidence. More recent work applies this idea to LLM-based systems: HMCF [cite] integrates GPT-4o with a human supervision interface for multi-robot rescue, allowing operators to issue natural-language corrections when the LLM's plan is unsatisfactory. Sirius [cite] provides a shared-control framework in which human corrections are used to adjust an autonomous agent's future behaviour. BrainBody-LLM [cite] uses closed-loop execution feedback—including human-in-the-loop signals—to trigger replanning in a single-robot manipulation task.

A common thread across all of these systems is that human corrections are communicated as natural language or low-level control signals, and are not structured at the protocol level. The LLM reads the correction and replans, but the system does not classify the correction semantically, does not validate it against a schema, and does not guarantee that the correction's intent will persist across subsequent planning steps. Our work addresses this gap directly.

\subsection{MCP in Robotic Systems}
MCP was introduced by Anthropic in late 2024 as a universal interface for connecting LLMs to external tools. Its adoption in robotics has accelerated rapidly: Robot MCP servers [cite] wrap ROS topics and services as MCP tools, allowing LLMs to command robots through natural language. IoT-MCP [cite] demonstrated MCP-based integration of sensor streams with LLM reasoning. MCP-Zero [cite] extends MCP with proactive tool discovery, allowing LLMs to request new capabilities on demand. To our knowledge, no prior work has used MCP as a channel for structured human authority injection—existing deployments treat MCP as a one-directional LLM-to-robot interface, not as a bidirectional human–LLM–robot coordination protocol.

\subsection{Positioning of This Work}

Table 1 positions our system relative to the closest prior work across five design dimensions. The combination of HITL override, MCP protocol layer, multi-robot rescue context, typed override classification, and persistent in-context propagation is unique to this paper.


\begin{table*}[t]
\centering
\caption{Comparison with closely related systems.}
\label{tab:comparison}

\small
\setlength{\tabcolsep}{4pt}
\renewcommand{\arraystretch}{1.1}

\begin{tabular}{lccccc}
\toprule
\textbf{System} & \textbf{HITL override} & \textbf{MCP layer} & \textbf{Multi-robot rescue} & \textbf{Structured override type} & \textbf{In-context propagation} \\
\midrule
HMCF [X] & Yes & No & Yes & No & No \\
BrainBody-LLM [X] & Yes & No & No & No & No \\
COHERENT [X] & No & No & Yes & No & No \\
Sirius [X] & Yes & No & No & No & No \\
This work & Yes & Yes & Yes & Yes & Yes \\
\bottomrule
\end{tabular}
\end{table*}